\documentclass[11pt]{article}

% ---- packages ----
\usepackage[margin=1in]{geometry}
\usepackage{amsmath,amssymb}
\usepackage{booktabs}
\usepackage{enumitem}
\usepackage{hyperref}
\usepackage{listings}
\usepackage{xcolor}

% ---- listings style for Python/Gurobi ----
\definecolor{codegray}{rgb}{0.95,0.95,0.95}
\definecolor{codekw}{rgb}{0.0,0.0,0.6}
\definecolor{codestr}{rgb}{0.58,0.0,0.0}
\definecolor{codecmt}{rgb}{0.25,0.5,0.35}
\lstdefinestyle{py}{
  backgroundcolor=\color{codegray},
  basicstyle=\ttfamily\footnotesize,
  keywordstyle=\color{codekw}\bfseries,
  stringstyle=\color{codestr},
  commentstyle=\color{codecmt}\itshape,
  showstringspaces=false,
  breaklines=true,
  frame=single,
  framesep=4pt,
  numbers=left,
  numberstyle=\tiny\color{gray},
  language=Python,
}

\newcommand{\R}{\mathbb{R}}
\newcommand{\Z}{\mathbb{Z}}

\title{\vspace{-1.5em}Assignment 4: Maximum Concurrent Flow and Network Design\\
\large CS 536 -- Data Communication and Computer Networks, Spring 2026}
\author{Pranay Nandkeolyar and Arav Popli}
\date{\today}

\begin{document}
\maketitle

\begin{center}
\small Code available at:
\url{https://github.com/PranayN23/cs536-assignment4}
\end{center}

\section*{Preliminaries}

We have a directed network $G$ with $n$ nodes and unit-capacity links,
where every node has exactly $d$ outgoing and $d$ incoming links.
Given a traffic matrix $T$ (with $T_{ii}=0$), the \emph{maximum concurrent
flow} $\alpha^*(G,T)$ is the largest factor $\alpha$ by which we can
scale $T$ and still route it without exceeding any link's capacity.
Intuitively, $\alpha^* = 0.7$ means we can deliver 70\% of the requested
traffic simultaneously. The goal is to design $G$ to make $\alpha^*$ as
large as possible.

\paragraph{Two upper bounds on $\alpha^*$.}
Both bounds come from a simple observation: the network has limited
capacity, so there is a ceiling on how much traffic can flow.

\textbf{Hop bound.} Since each node only has $d$ outgoing links, it can
only connect directly to $d$ neighbors; all other destinations require
multi-hop paths. The more hops traffic takes, the more link capacity it
consumes, which limits how much total demand can be routed. Summing
link capacities gives:
\begin{equation}
\alpha^*(G,T) \;\le\;
\frac{nd}{\displaystyle\sum_{s \neq t} T_{s,t}\,d_G(s,t)}.
\label{eq:hopbound}
\end{equation}

\textbf{Cut bound.} Any set of nodes $S$ can only send traffic to the
rest of the network through the links that cross the cut $(S, \bar{S})$:
\begin{equation}
\alpha^*(G,T) \;\le\;
\min_{S}
\frac{\text{links crossing cut }(S,\bar{S})}{\displaystyle\sum_{s\in S,\,t\notin S} T_{s,t}}.
\label{eq:cutbound}
\end{equation}

\bigskip\hrule\bigskip

\section*{Part 1 -- Uniform Traffic Matrix (50 pts)}

The uniform traffic matrix is
\[
T_{ij} = \frac{d}{n-1} \quad (i \neq j), \qquad T_{ii} = 0.
\]
Each row and column sums to exactly $d$, fully using each node's link
budget. This is the all-to-all pattern from AllGather/AllReduce steps
where every node communicates with every other node at the same rate.

\subsection*{Step 1: An upper bound from max-flow min-cut}

Think of the network as a shared pool of $nd$ unit-capacity link-slots
(each of the $n$ nodes contributes $d$ outgoing links). Routing demand
pair $(s,t)$ at rate $\alpha T_{st}$ consumes at least
$\alpha T_{st} \cdot d_G(s,t)$ link-slots, because the flow must
traverse at least $d_G(s,t)$ edges. Summing over all pairs and requiring
that total consumption cannot exceed total supply gives:
\[
\alpha \sum_{s \neq t} T_{st}\,d_G(s,t) \;\le\; nd.
\]
Rearranging and substituting $T_{st} = d/(n-1)$:
\[
\alpha^* \;\le\;
\frac{nd}{\dfrac{d}{n-1}\cdot n(n-1)\cdot\bar{h}(G)}
\;=\; \frac{1}{\bar{h}(G)},
\]
where $\bar{h}(G) = \frac{1}{n(n-1)}\sum_{s\neq t} d_G(s,t)$ is the
average hop count. \textbf{Maximizing $\alpha^*$ is therefore equivalent
to minimizing the average hop count.}

This is a max-flow min-cut argument in disguise: the ``cut'' is the
entire link budget of the network, and the bottleneck is how many times
each unit of demand must cross links on its way from source to
destination. Any topology that forces traffic to take long paths wastes
link capacity and reduces $\alpha^*$.

\subsection*{Step 2: A lower bound on average hop count (impossibility)}

How small can $\bar{h}(G)$ be? Since $\bar{h} \le D(G)$, it suffices to
lower-bound the diameter. In any $d$-regular network, a BFS from a
source discovers at most $d$ nodes at depth~1 and at most $d(d-1)^{k-1}$
new nodes at depth~$k \ge 2$ (each neighbor exposes at most $d-1$
previously unseen nodes). For all $n$ nodes to be reachable within
diameter $D$:
\[
n \;\le\; 1 + d\sum_{k=0}^{D-1}(d-1)^k \;\approx\; d \cdot (d-1)^{D-1},
\]
which forces $D \ge \log_{d-1} n$ (approximately). Hence
$\bar{h}(G) \ge \log_{d-1} n$ and:
\begin{equation}
\alpha^* \;\le\; \frac{1}{\log_{d-1} n}.
\label{eq:upper}
\end{equation}
No $d$-regular topology can beat this bound, no matter how cleverly
it is wired.

\subsection*{Step 3: The De~Bruijn digraph achieves the bound}

\paragraph{Construction.}
The De~Bruijn digraph $\mathrm{DB}(d,k)$ has $n = d^k$ nodes, each
represented as a length-$k$ string over $\{0,\ldots,d-1\}$. The edge
rule is: drop the leftmost character, append any new character on the
right. Formally, from node $b = (b_{k-1},\ldots,b_0)$ there is a
directed edge to each of the $d$ nodes
\[
b \;\longrightarrow\; (b_{k-2},\ldots,b_0,\,c) \quad \text{for each }
c \in \{0,\ldots,d-1\}.
\]
Every node has exactly $d$ outgoing and $d$ incoming edges, so the
graph is $d$-regular.

\begin{quote}\footnotesize
\textbf{Pseudocode: De~Bruijn digraph $\mathrm{DB}(d,k)$}
\begin{lstlisting}[language=,basicstyle=\ttfamily\footnotesize,frame=none,
                   numbers=none,backgroundcolor=\color{white}]
Input:  d (degree), k (string length), n = d^k
Output: directed d-regular graph on {0, ..., n-1}

for each node b in {0, ..., n-1}:
  prefix = b mod d^{k-1}          # drop the leading digit
  for c = 0, 1, ..., d-1:
    add edge  b -> prefix * d + c  # shift left, append c
\end{lstlisting}
\end{quote}

\paragraph{Routing.}
To route from $s$ to $t$, simply append the digits of $t$ one at a
time, most-significant first. After exactly $k$ steps, all digits of
$s$ have been shifted off and the current node is exactly $t$. Every
pair $(s,t)$ uses exactly $k = \log_d n$ hops---no path is longer or
shorter.

\paragraph{Why this is tight.}
Because every pair takes exactly $k$ hops, $\bar{h} = k$, and the
bound from Step~1 gives $\alpha^* \le 1/k$. To see that $\alpha = 1/k$
is achievable, route all demands simultaneously along their
$k$-hop De~Bruijn paths. The shift symmetry of the construction ensures
every link carries the same load: total traffic injected is $\alpha nd$,
each unit travels $k$ hops, so total link usage is $\alpha ndk$, spread
evenly across $nd$ links. Each link carries $\alpha k$; setting
$\alpha k = 1$ saturates every link simultaneously, so
$\alpha^* = 1/k = 1/\log_d n$, matching~\eqref{eq:upper} exactly.

\paragraph{Sanity check.}
When $d = n-1$ every node connects directly to every other node
($K_n$). Then $k=1$, $\bar{h}=1$, and $\alpha^*=1$---all traffic is
routed in one hop with no congestion, as expected.

\paragraph{What if $n$ is not a power of $d$?}
De~Bruijn requires $n = d^k$. For general $n$, a random $d$-regular
graph achieves an average hop count of approximately $\log_{d-1} n$,
giving the same asymptotic guarantee.

\paragraph{Why not a leaf-spine topology?}
A leaf-spine network achieves diameter~2 but requires the spine switches
to have degree~$n$, far exceeding the degree budget of $d$. If the
spine nodes are drawn from the same $n$-node pool (as required here,
since there are no additional switches), each spine can serve at most
$d$ hosts, collapsing to the complete graph $K_n$ when $d = n-1$. For
$d \ll n$, the degree constraint makes a true leaf-spine infeasible, and
De~Bruijn's $\log_d n$ diameter is the best achievable.

\paragraph{Summary.}
\emph{Use the De~Bruijn digraph. It minimizes average hop count to
$\log_d n$ via a simple shift-and-append construction, achieving
$\alpha^* = 1/\log_d n$, which matches the impossibility lower bound.
No $d$-regular topology can do better.}

\bigskip\hrule\bigskip

\section*{Part 2 -- Arbitrary (Hose) Traffic (50 pts)}

\paragraph{Motivation.}
Unlike Part~1, the traffic matrix here is not fixed in advance.
In practice---for example, in MoE inference where each token is routed
to a different expert GPU---the demand pattern is unpredictable and
changes every forward pass. We therefore want a topology that performs
well for \emph{any} $T$ in the hose polytope, not just a symmetric
all-to-all pattern. For a given $T$, we solve for the best topology
and routing jointly.

We fix $n=8$, $d=4$, and the hose polytope
\[
\mathcal{T}_{\text{hose}} = \big\{\, T \ge 0 : T_{ii}=0,\;
\textstyle\sum_j T_{ij} \le 4,\; \sum_i T_{ij} \le 4 \,\big\}.
\]
Given $T \in \mathcal{T}_{\text{hose}}$, we jointly choose topology $G$
and routing to maximize $\alpha^*(G,T)$.

\subsection*{MILP formulation}

We solve this as a Mixed-Integer Linear Program (MILP). The decision
variables are:
\begin{itemize}[leftmargin=*]
\item $x_{ij} \in \{0,1,\ldots,4\}$: number of directed links from node $i$ to node $j$ (the topology).
\item $f^{s,t}_{ij} \ge 0$: flow of demand pair $(s,t)$ on link $(i,j)$ (the routing).
\item $\alpha \ge 0$: the concurrent flow value to maximize.
\end{itemize}

\paragraph{Objective.} $\max\;\alpha$.

\paragraph{Constraints.}
\begin{align}
x_{ii} &= 0 && \text{(no self-loops)} \label{eq:noselfloop}\\
\sum_{j \neq i} x_{ij} &= 4 && \forall i \quad \text{(out-degree = 4)} \label{eq:outdeg}\\
\sum_{i \neq j} x_{ij} &= 4 && \forall j \quad \text{(in-degree = 4)} \label{eq:indeg}\\
\sum_{(s,t)} f^{s,t}_{ij} &\le x_{ij} && \forall i\neq j \quad \text{(link capacity)} \label{eq:cap}\\
\sum_{j} f^{s,t}_{vj} - \sum_{i} f^{s,t}_{iv}
  &= \alpha T_{s,t} \cdot \mathbf{1}[v{=}s] - \alpha T_{s,t} \cdot \mathbf{1}[v{=}t]
  && \forall (s,t), v \quad \text{(flow conservation)} \label{eq:flow}
\end{align}
Since $T_{s,t}$ is a given constant, the products $\alpha \cdot T_{s,t}$
are linear in $\alpha$, making this a valid LP/MILP.

\paragraph{What each constraint does.}
Constraints~\eqref{eq:outdeg}--\eqref{eq:indeg} ensure the chosen $x$
is a valid 4-regular topology. Constraint~\eqref{eq:cap} says total flow
on a link cannot exceed the number of parallel links wired there.
Constraint~\eqref{eq:flow} is standard flow conservation: for each
demand pair $(s,t)$, source $s$ injects $\alpha T_{s,t}$ and sink $t$
absorbs it, while all other nodes pass it through.

The MILP simultaneously decides:
\begin{enumerate}[nosep]
\item \textbf{Wiring}: how to allocate each node's 4 outgoing link slots
among the other 7 nodes.
\item \textbf{Routing}: how to split each demand $(s,t)$ over multi-hop
paths given that wiring.
\end{enumerate}

\subsection*{Implementation (Gurobi Python API)}

\lstinputlisting[style=py, firstline=50, lastline=165,
caption={\texttt{solution.py}: Gurobi MILP for Part~2 (core solver).}]{solution.py}

\paragraph{Implementation note: destination-aggregated flows.}
The per-pair formulation above has $O(n^2)$ flow variables per link.
To stay within Gurobi's free-license variable limit, we aggregate flows
by destination: let $g^t_{ij} = \sum_s f^{s,t}_{ij}$ be the total flow
on link $(i,j)$ heading toward destination $t$. The conservation
constraint becomes: every node $v \ne t$ sends out $\alpha T_{v,t}$ more
than it receives (it is a source for the $t$-flow), and $t$ absorbs the
rest. This reduces variables by a factor of $n$ and gives the same
optimal $\alpha$.

\subsection*{Results}

We tested three traffic matrices:
\begin{enumerate}[leftmargin=*,nosep]
\item \textbf{Cyclic permutation $\times 4$}: node $i$ sends all 4 units
to node $(i+1) \bmod 8$.
\item \textbf{Uniform hose}: $T_{ij} = 4/7$ for all $i \ne j$.
\item \textbf{Mixed}: two heavy demands plus a thin uniform background.
\end{enumerate}

\begin{center}
\begin{tabular}{lcc}
\toprule
Traffic matrix & $\alpha^*$ & Topology found \\
\midrule
Cyclic permutation $\times 4$ & $1.000$ & 4 parallel links $i \to (i{+}1)\bmod 8$ \\
Uniform hose $T_{ij}=4/7$     & $0.700$ & 4-regular digraph, all 32 links utilized \\
Mixed heavy + background      & $1.000$ & 3 direct links on heavy pairs, multi-hop background \\
\bottomrule
\end{tabular}
\end{center}

\paragraph{Checking the uniform-hose result.}
In any 4-regular graph on 8 nodes, each node has 4 neighbors at distance 1
and the remaining 3 nodes are at distance at least 2. So:
\[
\bar{h} \;\ge\; \frac{4 \cdot 1 + 3 \cdot 2}{7} = \frac{10}{7}.
\]
Plugging into the hop bound: $\alpha^* \le 1/\bar{h} \le 7/10 = 0.7$.
The MILP returns exactly $0.7$, so the topology it found is provably
optimal and every link is fully utilized.

\paragraph{Interpreting the results.}
\begin{itemize}[leftmargin=*]
\item When $T$ is an integer hose matrix (row/column sums $= 4$), direct
  routing is optimal: wire $x_{ij} = T_{ij}$ links and route one hop,
  giving $\alpha = 1$.
\item When $T$ is fractional (like $T_{ij}=4/7$), no integer wiring
  matches the demand directly. The MILP spreads links across multiple
  destinations and routes multi-hop, producing a well-connected topology
  consistent with the De~Bruijn intuition from Part~1.
\end{itemize}

\bigskip\hrule\bigskip

\section*{Deliverables}

\begin{itemize}[leftmargin=*,nosep]
\item \texttt{report.tex / report.pdf} -- this document.
\item \texttt{solution.py} -- Gurobi MILP implementation with three test
  cases and a constraint verifier.
\item \texttt{README.md} -- build and run instructions.
\item Full source at \url{https://github.com/PranayN23/cs536-assignment4}.
\end{itemize}

\end{document}